% !TEX root = ../glossario.tex
% ---------------------------------------------------------------------
%  Lettera A del Glossario. Una voce per comando, in ordine alfabetico:
%
%     \termine{Voce}{Definizione: una o due frasi chiare.}
%
%  Nella definizione si possono usare \emph{corsivo}, \textbf{grassetto},
%  \url{https://...} e \glo{altro termine} (se anche quello è nel
%  Glossario). Niente tabelle, immagini o elenchi.
%  Per una lettera nuova: copiate questo file con il nome della lettera
%  (es. capitoli/b.tex) e cambiate la riga \letteraglossario.
% ---------------------------------------------------------------------
\letteraglossario{A}

\termine{Analisi dei Requisiti}{Documento che descrive cosa deve fare il
  prodotto: casi d'uso e requisiti, ricavati dal capitolato e dagli
  incontri con il proponente.}

\termine{API}{\emph{Application Programming Interface}: insieme di
  funzioni che un programma mette a disposizione di altri programmi per
  usarne i servizi.}
